% Part: computability
% Chapter: computability-theory
% Section: applications-fixed-point

\documentclass[../../../include/open-logic-section]{subfiles}

\begin{document}

\olfileid[hi]{cmp}{thy}{apf}
\olsection{नियत-बिंदु प्रमेय का अनुप्रयोग}

नियत-बिंदु प्रमेय मूलतः हमें आंशिक संगणनीय फलनों को उनके सूचकांकों के
पदों में परिभाषित करने देता है। उदाहरणतः, हम ऐसा सूचकांक $e$ खोज सकते हैं
कि प्रत्येक $y$ के लिए
\[
\cfind{e}(y) = e + y.
\]
एक और उदाहरण के रूप में, नियत-बिंदु प्रमेय के प्रमाण का उपयोग Java या
C++ में ऐसा प्रोग्राम बनाने के लिए किया जा सकता है जो स्वयं को मुद्रित करे।

याद रखें कि यदि प्रत्येक $e$ के लिए $W_e$ को $\cfind{e}$ का प्रांत
मानें, तो अनुक्रम $W_0$, $W_1$, $W_2$,~\dots संगणनीयतः परिगणनीय
समुच्चयों का परिगणन करता है। इनमें से कुछ समुच्चय संगणनीय हैं। पूछा जा
सकता है कि क्या ऐसी कोई कलनविधि है जो मान $x$ को आगत ले और यदि $W_x$
संगणनीय हो, तो उसके अभिलाक्षणिक फलन का कोई सूचकांक लौटाए। उत्तर ``नहीं''
है; ऐसी कोई कलनविधि नहीं है:

\begin{thm}
निम्न गुण वाला कोई आंशिक संगणनीय फलन $f$ नहीं है: जब भी $W_e$ संगणनीय
हो, तब $f(e)$ परिभाषित हो और $\cfind{f(e)}$ उसका अभिलाक्षणिक फलन हो।
\end{thm}

\begin{proof}
$f$ को कोई संगणनीय फलन मानें; हम ऐसा $e$ बनाएँगे कि $W_e$ संगणनीय हो,
किंतु $\cfind{f(e)}$ उसका अभिलाक्षणिक फलन न हो। नियत-बिंदु प्रमेय से हम
ऐसा सूचकांक $e$ खोज सकते हैं कि
\[
\cfind{e}(y) \simeq
\begin{cases}
0 & \text{यदि $y=0$ और $\cfind{f(e)}(0) \fdefined = 0$} \\
\text{अपरिभाषित} & \text{अन्यथा।}
\end{cases}
\]
अर्थात् $e$ इस प्रकार परिभाषित फलन पर नियत-बिंदु प्रमेय लगाकर मिलता है:
\[
g(x,y) \simeq
\begin{cases}
0 & \text{यदि $y=0$ और $\cfind{f(x)}(0) \fdefined = 0$} \\
\text{अपरिभाषित} & \text{अन्यथा।}
\end{cases}
\]
अनौपचारिक रूप से हम इस प्रकार देख सकते हैं कि $g$ आंशिक संगणनीय है:
आगत $x$ और $y$ पर कलनविधि पहले जाँचती है कि $y$,~$0$ के बराबर है या
नहीं। यदि हाँ, तो वह $f(x)$ संगणित करती है और फिर सार्वत्रिक मशीन से
$\cfind{f(x)}(0)$ संगणित करती है। यदि यह अंतिम संगणना रुककर~$0$ लौटाती
है, तो कलनविधि~$0$ लौटाती है; अन्यथा वह नहीं रुकती।

अब ध्यान दें कि यदि $\cfind{f(e)}(0)$ परिभाषित और $0$ के बराबर है, तो
$\cfind{e}(y)$ ठीक तभी परिभाषित है जब $y$, $0$ के बराबर है; अतः
$W_e =
\{ 0 \}$। यदि $\cfind{f(e)}(0)$ अपरिभाषित है, अथवा परिभाषित
होकर भी $0$ के बराबर नहीं, तो $W_e = \emptyset$। दोनों ही स्थितियों में
$\cfind{f(e)}$,~$W_e$ का अभिलाक्षणिक फलन नहीं है, क्योंकि वह आगत $0$ पर
गलत उत्तर देता है।
\end{proof}

\end{document}
